\chapter{Architettura}
\label{cap:architettura}

Il servizio esiste in cinque versioni, tutte configurazioni dello stesso codice (D0012). Le versioni formano un albero di derivazione (figura~\ref{fig:derivazione}): ognuna cambia una sola caratteristica rispetto a quella da cui deriva, così ogni confronto isola un solo fattore.

\begin{figure}[htbp]
\centering
\begin{tikzpicture}
\node[box] (v1) {V1\\{\scriptsize sincrona, 1 istanza}};
\node[box, right=1.3cm of v1] (v2) {V2\\{\scriptsize N repliche}};
\node[box, above right=0.2cm and 1.6cm of v2] (v3) {V3\\{\scriptsize asincrona}};
\node[box, below right=0.2cm and 1.6cm of v2] (v4) {V4\\{\scriptsize DB replicato}};
\node[box, right=1.3cm of v3] (v5) {V5\\{\scriptsize autoscaling}};
\draw[freccia] (v1) -- node[above, font=\scriptsize] {+ repliche} (v2);
\draw[freccia] (v2) -- node[above, sloped, font=\scriptsize] {+ coda} (v3);
\draw[freccia] (v2) -- node[below, sloped, font=\scriptsize] {+ replica DB} (v4);
\draw[freccia] (v3) -- node[above, font=\scriptsize] {+ KEDA} (v5);
\end{tikzpicture}
\caption{Albero di derivazione delle versioni: ogni arco aggiunge un solo fattore.}
\label{fig:derivazione}
\end{figure}

\begin{table}[htbp]
\centering
\caption{Versioni del sistema.}
\label{tab:versioni}
\begin{tabularx}{\linewidth}{@{}llLl@{}}
\toprule
\textbf{Versione} & \textbf{Deriva da} & \textbf{Cosa cambia} & \textbf{Domande} \\
\midrule
V1 & --  & Riferimento: sincrona, una istanza dell'API & RQ1, RQ3 \\
V2 & V1  & N repliche dell'API dietro un bilanciatore & RQ1--RQ4, RQ6 \\
V3 & V2  & Comunicazione asincrona tramite coda e worker & RQ2, RQ5, RQ6 \\
V4 & V2  & Database replicato con failover automatico (V4a: replica asincrona, V4s: sincrona) & RQ4 \\
V5 & V3  & Numero di worker scalato sulla lunghezza della coda & RQ5 \\
\bottomrule
\end{tabularx}
\end{table}

\section{Componenti}

\begin{table}[htbp]
\centering
\caption{Componenti del sistema.}
\label{tab:componenti}
\begin{tabularx}{\linewidth}{@{}lllL@{}}
\toprule
\textbf{Componente} & \textbf{Stato} & \textbf{Piano} & \textbf{Responsabilità} \\
\midrule
API        & no & dati      & Riceve l'ordine; in V1, V2, V4 lo scrive nel database, in V3 e V5 lo pubblica in coda \\
Worker     & no & dati      & (V3, V5) Legge gli ordini dalla coda e li scrive nel database \\
RabbitMQ   & sì & dati      & (V3, V5) Coda degli ordini e coda dei messaggi falliti \\
PostgreSQL & sì & dati      & Archivio degli ordini; in V4 primario più standby \\
Toxiproxy  & no & dati      & Proxy verso il database, usato per iniettare latenza (F5); sempre presente in V2 e V3 \\
CloudNativePG & no & controllo & (V4) Rileva il guasto del primario e promuove lo standby \\
KEDA       & no & controllo & (V5) Legge la lunghezza della coda e regola il numero di worker \\
\bottomrule
\end{tabularx}
\end{table}

Il \textbf{piano dei dati} serve le richieste; il \textbf{piano di controllo} (Kubernetes, CloudNativePG, KEDA) le ripara e le adatta. Un guasto del piano di controllo non interrompe il servizio, ma impedisce il ripristino: per questo i due piani sono trattati separatamente nella sezione~\ref{sec:spof}.

\section{V1 -- sincrona, una istanza}

L'API riceve l'ordine, lo scrive nel database e risponde (figura~\ref{fig:v1}).

\begin{figure}[htbp]
\centering
\begin{tikzpicture}
\node[box] (client) {Client\\(JMeter)};
\node[box, right=1.6cm of client] (api) {API\\$\times 1$};
\node[box, right=1.6cm of api] (db) {PostgreSQL};
\draw[freccia] (client) -- (api);
\draw[freccia] (api) -- (db);
\end{tikzpicture}
\caption{V1: architettura sincrona con una sola istanza.}
\label{fig:v1}
\end{figure}

\section{V2 -- sincrona, N repliche}

Come V1, con N copie dell'API dietro un bilanciatore (figura~\ref{fig:v2}).

\begin{figure}[htbp]
\centering
\begin{tikzpicture}
\node[box] (client) {Client};
\node[box, right=1.3cm of client] (lb) {Service /\\Ingress};
\node[box, right=1.3cm of lb] (api) {API\\$\times N$};
\node[box, right=1.3cm of api] (db) {PostgreSQL};
\draw[freccia] (client) -- (lb);
\draw[freccia] (lb) -- (api);
\draw[freccia] (api) -- (db);
\end{tikzpicture}
\caption{V2: architettura sincrona con N repliche dell'API.}
\label{fig:v2}
\end{figure}

\section{V3 -- asincrona}

L'API pubblica l'ordine in coda, attende la conferma del broker e risponde ``accettato''. I worker leggono dalla coda, scrivono nel database e solo dopo confermano il messaggio. I messaggi che non si riesce a elaborare finiscono nella coda dei falliti (figura~\ref{fig:v3}).

\begin{figure}[htbp]
\centering
\resizebox{\linewidth}{!}{%
\begin{tikzpicture}
\node[box] (client) {Client};
\node[box, right=1.1cm of client] (lb) {Service /\\Ingress};
\node[box, right=1.1cm of lb] (api) {API\\$\times N$};
\node[box, right=1.6cm of api] (mq) {RabbitMQ};
\node[box, right=1.1cm of mq] (worker) {Worker\\$\times M$};
\node[box, right=1.1cm of worker] (db) {PostgreSQL};
\node[box, below=0.9cm of mq, minimum width=2cm] (dlq) {Dead-letter};
\draw[freccia] (client) -- (lb);
\draw[freccia] (lb) -- (api);
\draw[freccia] (api) -- node[above, font=\scriptsize] {publish} node[below, font=\scriptsize] {+ conferma} (mq);
\draw[freccia] (mq) -- (worker);
\draw[freccia] (worker) -- (db);
\draw[freccia, dashed] (mq) -- (dlq);
\end{tikzpicture}}
\caption{V3: architettura asincrona con coda e worker.}
\label{fig:v3}
\end{figure}

\section{V4 -- database replicato con failover}

Come V2, ma PostgreSQL è un cluster di due istanze gestito dall'operatore CloudNativePG: un \textbf{primario}, che accetta le scritture, e uno \textbf{standby}, che riceve in streaming il log delle transazioni (WAL). L'API si collega al servizio \texttt{-rw}, che punta sempre al primario corrente; se il primario cade, l'operatore promuove lo standby e aggiorna il servizio (figura~\ref{fig:v4}).

V4 si misura in due modalità, che rappresentano un compromesso classico tra durabilità, disponibilità e latenza:
\begin{description}
  \item[V4a -- replica asincrona] il commit non aspetta lo standby: latenza minima, ma gli ultimi commit possono andare persi nel failover.
  \item[V4s -- replica sincrona] il commit aspetta la conferma dello standby: nessuna perdita, latenza più alta, e le scritture si fermano se lo standby non risponde.
\end{description}

\begin{figure}[htbp]
\centering
\resizebox{\linewidth}{!}{%
\begin{tikzpicture}
\node[box] (client) {Client};
\node[box, right=1cm of client] (lb) {Service /\\Ingress};
\node[box, right=1cm of lb] (api) {API\\$\times N$};
\node[box, right=1cm of api] (rw) {Service\\\texttt{-rw}};
\node[box, right=1cm of rw] (prim) {PostgreSQL\\primario};
\node[box, below=1.1cm of prim] (stby) {PostgreSQL\\standby};
\node[box, right=1.4cm of prim, dashed] (op) {Operatore\\CloudNativePG};
\draw[freccia] (client) -- (lb);
\draw[freccia] (lb) -- (api);
\draw[freccia] (api) -- (rw);
\draw[freccia] (rw) -- (prim);
\draw[freccia] (prim) -- node[left, font=\scriptsize] {replica WAL} (stby);
\draw[freccia, dashed] (op) -- (prim);
\draw[freccia, dashed] (op) -- (stby);
\end{tikzpicture}}
\caption{V4: database primario e standby con failover gestito dall'operatore.}
\label{fig:v4}
\end{figure}

\section{V5 -- asincrona elastica}

Come V3, ma il numero di worker non è fisso: KEDA legge la lunghezza della coda e, tramite l'Horizontal Pod Autoscaler, porta i worker tra un minimo e un massimo (figura~\ref{fig:v5}). La reazione non è istantanea: si sommano l'intervallo di lettura della coda, la decisione dell'autoscaler e l'avvio dei pod. Il confronto con V3 si fa a parità di risorse consumate (D0015), altrimenti vincerebbe in modo banale la versione con più worker.

\begin{figure}[htbp]
\centering
\resizebox{\linewidth}{!}{%
\begin{tikzpicture}
\node[box] (client) {Client};
\node[box, right=1cm of client] (lb) {Service /\\Ingress};
\node[box, right=1cm of lb] (api) {API\\$\times N$};
\node[box, right=1.6cm of api] (mq) {RabbitMQ};
\node[box, right=1.4cm of mq] (worker) {Worker\\$1 \dots M$};
\node[box, right=1cm of worker] (db) {PostgreSQL};
\node[box, above=0.8cm of mq, minimum width=2cm] (dlq) {Dead-letter};
\coordinate (mid) at ($(mq)!0.5!(worker)$);
\node[box, below=1cm of mid, dashed] (keda) {KEDA};
\draw[freccia] (client) -- (lb);
\draw[freccia] (lb) -- (api);
\draw[freccia] (api) -- node[above, font=\scriptsize] {publish} (mq);
\draw[freccia] (mq) -- (worker);
\draw[freccia] (worker) -- (db);
\draw[freccia, dashed] (mq) -- (dlq);
\draw[freccia, dashed] (mq) -- node[left, font=\scriptsize] {lunghezza coda} (keda);
\draw[freccia, dashed] (keda) -- node[right, font=\scriptsize] {scala} (worker);
\end{tikzpicture}}
\caption{V5: worker scalati da KEDA sulla lunghezza della coda.}
\label{fig:v5}
\end{figure}

\section{Punti unici di guasto}
\label{sec:spof}

\begin{table}[htbp]
\centering
\caption{Punti unici di guasto per versione.}
\label{tab:spof}
\begin{tabularx}{\linewidth}{@{}lLL@{}}
\toprule
\textbf{Versione} & \textbf{Piano dei dati} & \textbf{Piano di controllo} \\
\midrule
V1 & Punto di ingresso, API, PostgreSQL & Control plane di Kubernetes (un nodo) \\
V2 & Punto di ingresso, PostgreSQL & Control plane \\
V3 & Punto di ingresso, RabbitMQ (accettazione), PostgreSQL (completamento) & Control plane \\
V4 & Punto di ingresso & Control plane, operatore (senza operatore niente failover) \\
V5 & Come V3 & Control plane, KEDA (senza KEDA i worker restano al numero corrente) \\
\bottomrule
\end{tabularx}
\end{table}

Il punto di ingresso è un limite dell'ambiente locale: è la porta di un nodo esposta all'host, che nessun esperimento guasta. Su AWS è l'ALB, distribuito su più zone.
